% ─────────────────────────────────────────────────────────────────────────────
\subsection*{Overview}

The Networked Control System (NCS) from \textbf{CKW10~§7.3} is a hybrid
continuous/discrete closed-loop control system.
A second-order continuous plant is integrated via QSS1; a STDEVS sensor
adds AWGN and samples periodically; a STDEVS network delay (NIRD) introduces
a uniformly distributed transmission delay into the feedback path.
The experiment sweeps the maximum delay $\tau_{\max}$ and reports the
mean-square tracking error MSE~$= \tfrac{1}{T}\int_0^T (1-y(t))^2\,dt$,
which increases with delay and diverges as $\tau_{\max} \to h$.

% ─────────────────────────────────────────────────────────────────────────────
\subsection*{System Topology}

\begin{center}
\begin{tikzpicture}[
  box/.style={draw,rectangle,minimum width=2.2cm,minimum height=0.85cm,
              align=center,font=\small},
  >=Stealth, node distance=0.55cm and 1.6cm
]
  \node[box] (plant) {Plant\\$G_P(s)$};
  \node[box, right=of plant] (sensor) {Sensor\\(ZOH+AWGN)};
  \node[box, right=of sensor] (nird) {NIRD};
  \node[box, below=1.0cm of nird] (hc) {Hold+Ctrl\\$u=1-y_h$};
  \node[box, below=1.0cm of plant] (wsum) {wsum\\$\dot{x}_2=u-0.8x_2$};

  \draw[->] (plant) -- node[above,font=\scriptsize]{$y(t)$} (sensor);
  \draw[->] (sensor) -- node[above,font=\scriptsize]{$y'(kh)$} (nird);
  \draw[->] (nird) -- node[right,font=\scriptsize]{$y'(kh{+}\tau)$} (hc);
  \draw[->] (hc) -- node[below,font=\scriptsize]{$u$} (wsum);
  \draw[->] (wsum) -- node[left,font=\scriptsize]{$\dot{x}_2$} (plant);
\end{tikzpicture}
\end{center}

The plant $G_P(s) = 1/(s^2 + 0.8s)$ is realised in state-space form:
$\dot{x}_1 = x_2$ and $\dot{x}_2 = u - 0.8\,x_2$.
Each integrator is a QSS1 atom.
The wsum block computes $\dot{x}_2 = 1\cdot u + (-0.8)\cdot x_2$; it also
receives $x_2$ from the second integrator for the feedback term.

\paragraph{Internal couplings (QSS subsystem).}
\begin{align*}
  \mathrm{int}_2.\texttt{out\_q} &\to \mathrm{int}_1.\texttt{in\_u}
    & (x_2 = \dot{x}_1) \\
  \mathrm{int}_2.\texttt{out\_q} &\to \mathrm{wsum}.\texttt{in}\langle 1\rangle
    & (x_2\;\text{for}\;{-0.8}\,x_2\;\text{term}) \\
  \mathrm{wsum}.\texttt{out}    &\to \mathrm{int}_2.\texttt{in\_u}
    & (\dot{x}_2) \\
  \mathrm{int}_1.\texttt{out\_q} &\to \mathrm{sensor}.\texttt{in\_y}
    & (\text{ZOH update of}\;y)
\end{align*}

\paragraph{Stochastic/control path.}
\begin{align*}
  \mathrm{sensor}.\texttt{out\_y}    &\to \mathrm{nird}.\texttt{in\_y} \\
  \mathrm{nird}.\texttt{out\_y}      &\to \mathrm{hold\_ctrl}.\texttt{in\_y} \\
  \mathrm{hold\_ctrl}.\texttt{out\_u}&\to \mathrm{wsum}.\texttt{in}\langle 0\rangle
\end{align*}

% ─────────────────────────────────────────────────────────────────────────────
\subsection*{Atomic Model Specifications}

\paragraph{Plant integrators (deterministic DEVS / QSS1).}

Two instances of the QSS1 integrator from \texttt{cadmium::basic\_models::qss},
with absolute quantum $\Delta q = 0.01$.  The first integrates $\dot{x}_1 = x_2$
(output is $y$); the second integrates $\dot{x}_2 = u - 0.8\,x_2$.

\paragraph{Sensor (STDEVS atomic).}

Combines the Sense+Sample and AWGN Generator components from the paper into
a single atom to avoid the SELECT ambiguity that would arise if both fired
at the same instant~$kh$.

\[
  M^{\mathrm{sensor}} = \langle X, Y, S, \delta_{\mathrm{int}},
    \delta_{\mathrm{ext}}, \lambda, ta \rangle
\]
\begin{align*}
  X &= \mathbb{R}\times\{\texttt{in\_y}\},\quad
  Y  = \mathbb{R}\times\{\texttt{out\_y}\} \\
  S &= (y_h \in \mathbb{R},\;\eta \in \mathbb{R},\;\sigma \in \mathbb{R}^+)\\
  \delta_{\mathrm{ext}}\bigl((y_h,\eta,\sigma),\,e,\,x\bigr)
    &= (x,\;\eta,\;\sigma - e) \\
  \delta_{\mathrm{int}}\bigl((y_h,\eta,\sigma),\,r\bigr)
    &= \bigl(y_h,\;\tilde\eta,\;h\bigr), \\
  &\quad \tilde\eta
       = \sqrt{-2\ln r_1}\cos(2\pi r_2)\sqrt{v_\eta},
       \quad (r_1,r_2)\sim U^2(0,1) \\
  \lambda(y_h,\eta,\sigma) &= y_h + \eta \\
  ta(y_h,\eta,\sigma) &= \sigma
\end{align*}

External inputs update the held plant output $y_h$ (ZOH); the new
noise sample $\tilde\eta$ is drawn at each internal transition using the
Box--M\"uller transform with \texttt{std::mt19937}.

\paragraph{NIRD (STDEVS atomic).}

\[
  M^{\mathrm{NIRD}} = \langle X, Y, S, \delta_{\mathrm{int}},
    \delta_{\mathrm{ext}}, \lambda, ta \rangle
\]
\begin{align*}
  X &= \mathbb{R}\times\{\texttt{in\_y}\},\quad
  Y  = \mathbb{R}\times\{\texttt{out\_y}\} \\
  S &= (y' \in \mathbb{R},\;\sigma \in \mathbb{R}^+_0\cup\{+\infty\}) \\
  \delta_{\mathrm{ext}}\bigl((y',\sigma),\,e,\,x,\,r\bigr)
    &= \begin{cases}
         (x,\;r\cdot\tau_{\max}) & \sigma = +\infty\;\text{(passive)} \\
         (y',\;\sigma - e)       & \text{otherwise (drop arrival)}
       \end{cases} \\
  \delta_{\mathrm{int}}\bigl((y',\sigma),\,r\bigr) &= (y',\;+\infty) \\
  \lambda(y',\sigma) &= y',\quad ta(y',\sigma) = \sigma
\end{align*}

Single-buffer: a new sample arriving while a previous one is in transit is
discarded.

\paragraph{Hold + Controller (deterministic DEVS).}

Receives the delayed noisy sample $y'(kh+\tau)$, holds it, and outputs
$u = 1 - y_h$.  Fires at $t=0$ with the initial control signal $u=1$
(reference minus initial plant output $y(0)=0$).

% ─────────────────────────────────────────────────────────────────────────────
\subsection*{Design Decisions}

\begin{enumerate}

\item \textbf{Manual event loop.}
  The QSS integrators are plain DEVS models whose transitions take no RNG,
  while the sensor and NIRD are STDEVS models.
  The cadmium STDEVS runner requires all models to accept a URNG; mixing
  both kinds in a single coupled model would therefore require wrapping the
  QSS atoms.  A manual event loop avoids this and follows the same pattern
  as the LBM experiment.

\item \textbf{Elapsed-time update discipline.}
  In the DEVS formalism, the elapsed time $e$ passed to
  $\delta_{\mathrm{ext}}$ is the time since the model's \emph{last
  transition}---internal or external.  The \texttt{last\_*} timestamps must
  therefore be updated after every internal transition, not just after
  external ones.  Omitting this causes the QSS integrators to double-advance
  their state variable, producing near-Zeno event cascades and incorrect
  trajectories.

\item \textbf{Sensor merges AWGN and Sense+Sample.}
  In the paper, the AWGN generator and Sense+Sample are separate atoms that
  both fire at~$kh$.  Classic DEVS SELECT would pick one arbitrarily;
  merging them into a single atom sidesteps the tie and is consistent with
  the paper's intent (noise and sample happen at the same instant).

\item \textbf{Quantum.}
  $\Delta q = 0.01$ provides sufficient accuracy to reproduce the
  qualitative MSE trend; $\Delta q = 0.001$ would add roughly $10\times$
  more events with no meaningful improvement in the MSE comparison.

\end{enumerate}

% ─────────────────────────────────────────────────────────────────────────────
\subsection*{Parameters}

\begin{center}
\begin{tabular}{lll}
\hline
Parameter & Symbol & Value \\
\hline
Sampling period & $h$ & $1$~s \\
AWGN variance & $v_\eta$ & $0.001$ \\
QSS quantum & $\Delta q$ & $0.01$ \\
Network delay range & $\tau_{\max}$ & $0.1,\,0.2,\ldots,0.8$~s \\
Simulation time & $T$ & $10\,000$~s \\
Repetitions & $N$ & $15$ \\
RNG & & \texttt{std::mt19937}, seed $42+\text{rep}$ \\
\hline
\end{tabular}
\end{center}
